\textbf{COLUNI 2023 -- Questão 29}\par\medskip Em uma série de exames médicos de rotina, que incluía exames de sangue, urina e fezes, D. Maria foi diagnosticada com diabetes. Muito preocupada, D. Maria relatou o fato ao marido, que respondeu: “Ih, isso é problema no sangue!”. Quando levou os resultados ao seu médico, ele explicou a ela que o exame de sangue detecta essa e várias outras doenças e que a diabetes, na verdade, pode ocorrer devido ao mau funcionamento de um determinado órgão.

\par

A alternativa que apresenta CORRETAMENTE o órgão ao qual o médico se refere é:\begin{enumerate}
\item pâncreas, responsável pela produção do hormônio insulina.
\item glândulas salivares, responsáveis pela digestão inicial de açúcares na boca.
\item baço, responsável pela destruição de células sanguíneas velhas.
\item rins, responsáveis pela filtração do sangue e eliminação de excretas.\end{enumerate}
